\documentclass[10pt,a4paper]{amsart}
\usepackage{amsmath,amssymb,amsthm,graphicx}
\usepackage{verbatim}
\newenvironment{reference}{\comment}{\endcomment}
%\newenvironment{reference}{}{}
\newenvironment{slogan}{\comment}{\endcomment}
\newenvironment{history}{\comment}{\endcomment}
% For commutative diagrams we use Xy-pic
\usepackage[all]{xy}

% We use 2cell for 2-commutative diagrams.
\xyoption{2cell}
% Theorem environments.
%
\theoremstyle{plain}
\newtheorem{theorem}[subsection]{Theorem}
\newtheorem{proposition}[subsection]{Proposition}
\newtheorem{lemma}[subsection]{Lemma}

\theoremstyle{definition}
\newtheorem{definition}[subsection]{Definition}
\newtheorem{example}[subsection]{Example}
\newtheorem{exercise}[subsection]{Exercise}
\newtheorem{situation}[subsection]{Situation}

\theoremstyle{remark}
\newtheorem{remark}[subsection]{Remark}
\newtheorem{remarks}[subsection]{Remarks}

\numberwithin{equation}{subsection}

% Macros
%
\def\lim{\mathop{\mathrm{lim}}\nolimits}
\def\colim{\mathop{\mathrm{colim}}\nolimits}
\def\Spec{\mathop{\mathrm{Spec}}}
\def\Hom{\mathop{\mathrm{Hom}}\nolimits}
\def\Ext{\mathop{\mathrm{Ext}}\nolimits}
\def\SheafHom{\mathop{\mathcal{H}\!\mathit{om}}\nolimits}
\def\SheafExt{\mathop{\mathcal{E}\!\mathit{xt}}\nolimits}
\def\Sch{\mathit{Sch}}
\def\Mor{\mathop{\mathrm{Mor}}\nolimits}
\def\Ob{\mathop{\mathrm{Ob}}\nolimits}
\def\Sh{\mathop{\mathit{Sh}}\nolimits}
\def\NL{\mathop{N\!L}\nolimits}
\def\CH{\mathop{\mathrm{CH}}\nolimits}
\def\proetale{{pro\text{-}\acute{e}tale}}
\def\etale{{\acute{e}tale}}
\def\QCoh{\mathit{QCoh}}
\def\Ker{\mathop{\mathrm{Ker}}}
\def\Im{\mathop{\mathrm{Im}}}
\def\Coker{\mathop{\mathrm{Coker}}}
\def\Coim{\mathop{\mathrm{Coim}}}

% Boxtimes
%
\DeclareMathSymbol{\boxtimes}{\mathbin}{AMSa}{"02}

%
% Macros for moduli stacks/spaces
%
\def\QCohstack{\mathcal{QC}\!\mathit{oh}}
\def\Cohstack{\mathcal{C}\!\mathit{oh}}
\def\Spacesstack{\mathcal{S}\!\mathit{paces}}
\def\Quotfunctor{\mathrm{Quot}}
\def\Hilbfunctor{\mathrm{Hilb}}
\def\Curvesstack{\mathcal{C}\!\mathit{urves}}
\def\Polarizedstack{\mathcal{P}\!\mathit{olarized}}
\def\Complexesstack{\mathcal{C}\!\mathit{omplexes}}
% \Pic is the operator that assigns to X its picard group, usage \Pic(X)
% \Picardstack_{X/B} denotes the Picard stack of X over B
% \Picardfunctor_{X/B} denotes the Picard functor of X over B
\def\Pic{\mathop{\mathrm{Pic}}\nolimits}
\def\Picardstack{\mathcal{P}\!\mathit{ic}}
\def\Picardfunctor{\mathrm{Pic}}
\def\Deformationcategory{\mathcal{D}\!\mathit{ef}}
\usepackage{cleveref}
\crefname{originaldefinition}{original definition}{original definitions}
\Crefname{originaldefinition}{Original definition}{Original definitions}
\makeatletter
\newlabel{definition-quasi-compact}{{definition-quasi-compact}{source}{}{sourcepointer-1}{}}
\newlabel{definition-quasi-compact@cref}{{[originaldefinition][0][]definition-quasi-compact}{[1][0][]source}{}{}{}}
\makeatother
\crefname{originallemma}{original lemma}{original lemmas}
\Crefname{originallemma}{Original lemma}{Original lemmas}
\makeatletter
\newlabel{lemma-base-change-quasi-compact}{{lemma-base-change-quasi-compact}{source}{}{sourcepointer-2}{}}
\newlabel{lemma-base-change-quasi-compact@cref}{{[originallemma][0][]lemma-base-change-quasi-compact}{[1][0][]source}{}{}{}}
\makeatother
\crefname{originalreference}{original reference}{original references}
\Crefname{originalreference}{Original reference}{Original references}
\makeatletter
\newlabel{morphisms-lemma-reach-points-scheme-theoretic-image}{{morphisms-lemma-reach-points-scheme-theoretic-image}{source}{}{sourcepointer-3}{}}
\newlabel{morphisms-lemma-reach-points-scheme-theoretic-image@cref}{{[originalreference][0][]morphisms-lemma-reach-points-scheme-theoretic-image}{[1][0][]source}{}{}{}}
\makeatother
\crefname{originalreference}{original reference}{original references}
\Crefname{originalreference}{Original reference}{Original references}
\makeatletter
\newlabel{stacks-properties-lemma-points-cartesian}{{stacks-properties-lemma-points-cartesian}{source}{}{sourcepointer-4}{}}
\newlabel{stacks-properties-lemma-points-cartesian@cref}{{[originalreference][0][]stacks-properties-lemma-points-cartesian}{[1][0][]source}{}{}{}}
\makeatother
\crefname{originalreference}{original reference}{original references}
\Crefname{originalreference}{Original reference}{Original references}
\makeatletter
\newlabel{stacks-properties-lemma-topology-points}{{stacks-properties-lemma-topology-points}{source}{}{sourcepointer-5}{}}
\newlabel{stacks-properties-lemma-topology-points@cref}{{[originalreference][0][]stacks-properties-lemma-topology-points}{[1][0][]source}{}{}{}}
\makeatother
\title{A point of the target of a quasi-compact morphism reached by the image lies in the scheme theoretic image}
\author{The Stacks Project Authors}
\date{stacks-project:stacks-morphisms.tex:local-snapshot-20261012; GNU FDL 1.2 or later}
\begin{document}
\maketitle
\noindent\textit{Editorial note.} All delivered mathematical text is the author's own, copied from the pinned Stacks Project chapter source sources/stacks-algebraic-stacks-morphisms/pinned-original-source.tex; no statement or proof step was rewritten, reordered or removed.\par\medskip
\noindent\textit{Editorial note.} Preamble provenance: the Stacks Project keeps one single shared preamble for every chapter; this card takes it directly from the staged sources/stacks-algebraic-stacks-morphisms/preamble.tex as three verbatim spans (lines 8-12, 14-18 and 154-217), so no derived file was written and nothing was added or redefined. The wrapper class and the base packages are supplied by the assembler.\par\medskip
\noindent\textit{Editorial note.} Bibliography provenance: neither the counted statement nor its proof contains a citation (this lane rescans the whole span before writing the card), so the card carries no bibliography and no entry was reconstructed.\par\medskip
\noindent\textit{Editorial note.} Selection provenance: this card was opened under the file-based control contract (CONTROL.md), with the mechanical pre-card gate run on its statement and proof interval before the card was written.\par\medskip
\noindent\textit{Allowed standard prerequisites (editorial).}\begin{itemize}
\item The description of points reached by the image of a morphism of algebraic stacks. Invoked inside the counted proof (source lines 2116-2125). Reformulates the hypothesis as a specialization condition on points.
\item The reduction of the scheme theoretic image along a flat base change and along a smooth cover of the source. Invoked inside the counted proof (source lines 2126-2140). Carries the statement to a morphism of schemes.
\item The corresponding statement for schemes that a point in the closure of the image lies in the scheme theoretic image. Invoked at the end of the counted proof (source lines 2141-2149) as a cross-chapter pointer. Supplies the conclusion after the reduction.
\end{itemize}\medskip
\ifdefined\hypertarget\hypertarget{sourcepointer-1}{}\fi
\noindent\textit{Source reference.} Original definition definition-quasi-compact: sources/stacks-algebraic-stacks-morphisms/pinned-original-source.tex, source lines 1967--1967. Navigation pointer only: the counted text refers to the Stacks item definition-quasi-compact, which is an existing labeled result of this same chapter and is not part of this card. The pointer carries the target's own label as its number and its definition line as its span; its content is not reproduced here.\par\medskip
\ifdefined\hypertarget\hypertarget{sourcepointer-2}{}\fi
\noindent\textit{Source reference.} Original lemma lemma-base-change-quasi-compact: sources/stacks-algebraic-stacks-morphisms/pinned-original-source.tex, source lines 1984--1984. Navigation pointer only: the counted text refers to the Stacks item lemma-base-change-quasi-compact, which is an existing labeled result of this same chapter and is not part of this card. The pointer carries the target's own label as its number and its definition line as its span; its content is not reproduced here.\par\medskip
\ifdefined\hypertarget\hypertarget{sourcepointer-3}{}\fi
\noindent\textit{Source reference.} Original reference morphisms-lemma-reach-points-scheme-theoretic-image: sources/stacks-algebraic-stacks-morphisms/pinned-original-source.tex, source lines 2141--2141. Navigation pointer only: the counted text refers to the Stacks item morphisms-lemma-reach-points-scheme-theoretic-image, which belongs to another chapter of the project and is not part of this card. The pointer carries the target's own label as its number because this chapter's staged snapshot does not contain the other chapter's tag table; its content is not reproduced here.\par\medskip
\ifdefined\hypertarget\hypertarget{sourcepointer-4}{}\fi
\noindent\textit{Source reference.} Original reference stacks-properties-lemma-points-cartesian: sources/stacks-algebraic-stacks-morphisms/pinned-original-source.tex, source lines 2126--2126. Navigation pointer only: the counted text refers to the Stacks item stacks-properties-lemma-points-cartesian, which belongs to another chapter of the project and is not part of this card. The pointer carries the target's own label as its number because this chapter's staged snapshot does not contain the other chapter's tag table; its content is not reproduced here.\par\medskip
\ifdefined\hypertarget\hypertarget{sourcepointer-5}{}\fi
\noindent\textit{Source reference.} Original reference stacks-properties-lemma-topology-points: sources/stacks-algebraic-stacks-morphisms/pinned-original-source.tex, source lines 2128--2128. Navigation pointer only: the counted text refers to the Stacks item stacks-properties-lemma-topology-points, which belongs to another chapter of the project and is not part of this card. The pointer carries the target's own label as its number because this chapter's staged snapshot does not contain the other chapter's tag table; its content is not reproduced here.\par\medskip
\begin{lemma}
\label{lemma-reach-points-scheme-theoretic-image}
Let $f : \mathcal{X} \to \mathcal{Y}$ be a quasi-compact morphism of
algebraic stacks. Let $y \in |\mathcal{Y}|$ be a point in the closure
of the image of $|f|$. There exists a valuation ring $A$ with
fraction field $K$ and a commutative diagram
$$
\xymatrix{
\Spec(K) \ar[r] \ar[d] & \mathcal{X} \ar[d] \\
\Spec(A) \ar[r] & \mathcal{Y}
}
$$
such that the closed point of $\Spec(A)$ maps to $y$.
\end{lemma}

\begin{proof}
Choose an affine scheme $V$ and a point $v \in V$ and a smooth morphism
$V \to \mathcal{Y}$ sending $v$ to $y$. Consider the base change diagram
$$
\xymatrix{
V \times_\mathcal{Y} \mathcal{X} \ar[r] \ar[d]_g & \mathcal{X} \ar[d]^f \\
V \ar[r] & \mathcal{Y}
}
$$
Recall that $|V \times_\mathcal{Y} \mathcal{X}| \to
|V| \times_{|\mathcal{Y}|} |\mathcal{X}|$ is surjective
(Properties of Stacks, Lemma \ref{stacks-properties-lemma-points-cartesian}).
Because $|V| \to |\mathcal{Y}|$ is open
(Properties of Stacks, Lemma \ref{stacks-properties-lemma-topology-points})
we conclude that $v$ is in the closure of the image of $|g|$.
Thus it suffices to prove the lemma for the quasi-compact morphism $g$
(Lemma \ref{lemma-base-change-quasi-compact}) which we do in the next
paragraph.

\medskip\noindent
Assume $\mathcal{Y} = Y$ is an affine scheme. Then $\mathcal{X}$
is quasi-compact as $f$ is quasi-compact
(Definition \ref{definition-quasi-compact}).
Choose an affine scheme $W$ and a surjective smooth morphism
$W \to \mathcal{X}$. Then the image of $|f|$ is the image
of $W \to Y$.
By Morphisms, Lemma \ref{morphisms-lemma-reach-points-scheme-theoretic-image}
we can choose a diagram
$$
\xymatrix{
\Spec(K) \ar[r] \ar[d] & W \ar[d] \ar[r] & \mathcal{X} \ar[d] \\
\Spec(A) \ar[r] & Y \ar[r] & Y
}
$$
such that the closed point of $\Spec(A)$ maps to $y$.
Composing with $W \to \mathcal{X}$ we obtain a solution.
\end{proof}

\end{document}
